\chapter{Resultados e Análise Computacional} \label{ch:resultados}

Este capítulo apresenta os experimentos computacionais realizados para calibrar e avaliar a matheurística LNS proposta. Inicialmente, a Seção~\ref{sec:instancias} descreve as instâncias de teste, o ambiente computacional e a metodologia de calibração adotada. Em seguida, a Seção~\ref{sec:analise_interna} detalha a análise incremental dos três parâmetros da LNS --- taxa de destruição, sentido de ordenação dos candidatos e tempo máximo do sub-MIP ---, conduzida conforme a abordagem \emph{one-factor-at-a-time}. Por fim, a Seção~\ref{sec:comparativo} apresenta o comparativo de desempenho entre a combinação de parâmetros selecionada para a LNS, o B\&B puro, a heurística GRASP e a variante híbrida GRASP|LNS nas 42 instâncias sintéticas e em uma instância real, avaliando tempo de execução, gap de otimalidade, número de nós restantes na árvore de busca e execuções efetivas da LNS.

\section{Instâncias e Ambiente de Testes} \label{sec:instancias}

Os experimentos computacionais foram conduzidos utilizando um conjunto de 42 instâncias de teste geradas sinteticamente, abrangendo diferentes cenários do Problema da Distribuição de Disciplinas (PDD). As instâncias variam em tamanho (62 a 214 disciplinas), número de professores (22 a 69), complexidade das restrições de carga horária e distribuição de preferências entre áreas de conhecimento. Para o comparativo final, acrescentou-se uma instância real, identificada nas tabelas como \texttt{entrada\_real}, composta por 224 disciplinas e 61 professores. Dessa forma, a avaliação geral considera 43 instâncias, mantendo a instância real distinguível das sintéticas e, simultaneamente, incorporada às métricas agregadas.

As instâncias foram produzidas pelo procedimento sintetizado no Algoritmo~\ref{alg:geracao_instancias}. A quantidade de tipos de disciplina foi fornecida como parâmetro e assumiu os valores 12, 16, 20, 25 e 32 nas instâncias utilizadas. Um tipo de disciplina representa uma categoria-base, como Cálculo ou Programação, e define uma área de conhecimento. Para cada tipo, o gerador cria entre 1 e 10 ofertas; cada oferta corresponde a uma disciplina individual do PDD, com carga horária própria, que deverá ser atribuída a um professor.

O fator de limite de professores foi mantido igual a 2 e utilizado somente para calcular a quantidade máxima de docentes que poderia ser gerada em cada área, dada por $\max\{1,\lfloor\text{número de ofertas}/2\rfloor\}$. A quantidade efetiva de professores é sorteada entre 1 e esse limite; portanto, o fator não estabelece uma proporção exata entre disciplinas e professores. Em seguida, são gerados os limites de carga horária e as preferências docentes. O ajuste final da carga máxima assegura que a capacidade docente da área seja suficiente para cobrir a carga horária de suas ofertas.

\begingroup
\renewcommand{\lstlistingname}{Algoritmo}
\begin{lstlisting}[
    caption={Pseudocódigo para geração das instâncias sintéticas.},
    label={alg:geracao_instancias},
    basicstyle=\ttfamily\footnotesize,
    frame=single,
    breaklines=true,
    numbers=none,
    columns=fullflexible,
    keepspaces=true,
    showstringspaces=false,
    float=htbp
]
ENTRADA:
    numero_tipos <- valor de entrada em {12, 16, 20, 25, 32}
    fator_limite_professores <- 2

PARA cada tipo de disciplina FACA
    criar uma area de conhecimento exclusiva para o tipo
    numero_ofertas <- inteiro aleatorio entre 1 e 10

    PARA cada oferta do tipo FACA
        carga_horaria <- inteiro aleatorio entre 0 e 100
        registrar a oferta no primeiro semestre e na area do tipo
    FIM PARA

    carga_total_area <- soma das cargas horarias das ofertas
    maximo_professores <- max(1, piso(numero_ofertas /
                                      fator_limite_professores))
    numero_professores <- inteiro aleatorio entre 1 e
                           maximo_professores
    carga_media <- piso(carga_total_area / numero_professores)

    PARA cada professor da area FACA
        sortear carga minima entre 0 e carga_media
        sortear carga maxima entre a carga minima e carga_total_area
        numero_preferencias <- inteiro aleatorio entre 1 e
                                numero_ofertas
        selecionar, sem repeticao, esse numero de ofertas da area

        PARA cada oferta selecionada FACA
            peso_preferencia <- inteiro aleatorio entre 0 e 10
        FIM PARA
    FIM PARA

    SE a soma das cargas maximas nao superar carga_total_area ENTAO
        aumentar a carga maxima do ultimo professor ate que a soma
        supere carga_total_area
    FIM SE
FIM PARA

gravar totais, ofertas, professores e preferencias em arquivo CSV
\end{lstlisting}
\endgroup

Na calibração interna apresentada na Seção~\ref{sec:analise_interna}, foram selecionadas 10 instâncias dentre as 42 sintéticas disponíveis em cada etapa, mantendo o tempo total de experimentação viável. As mesmas instâncias foram utilizadas entre os cenários comparados dentro de cada etapa, assegurando comparabilidade direta dos parâmetros. Após a calibração, o comparativo final da Seção~\ref{sec:comparativo} foi reexecutado com todas as 42 instâncias sintéticas e com a instância real.

Todas as execuções foram realizadas com limite de tempo global de 400 segundos por instância, em uma máquina equipada com processador AMD Ryzen 5 5600G, 64\,GB de memória RAM e sistema operacional Arch Linux. O solver utilizado foi o SCIP 7.0 \cite{gamrath_2020_scip}, integrado com as heurísticas LNS e GRASP desenvolvidas neste trabalho. A calibração dos parâmetros da heurística LNS foi conduzida de forma incremental, conforme a abordagem \emph{one-factor-at-a-time} (OFAT) \cite{montgomery_2017}: primeiro determinou-se a melhor taxa de destruição, em seguida o melhor sentido de ordenação e, finalmente, o melhor limite de tempo para o sub-MIP. Essa estratégia ajusta um parâmetro por vez, mantendo os demais fixados, o que permite isolar o efeito individual de cada fator e identificar a configuração mais adequada de forma sistemática. A escolha pelo OFAT foi motivada principalmente pelo custo computacional de uma calibração exaustiva: um projeto fatorial completo sobre os três parâmetros (4 níveis de taxa de destruição $\times$ 2 sentidos de ordenação $\times$ 4 níveis de tempo do sub-MIP = 32 combinações), com 10 instâncias e limite de 400 segundos cada, demandaria até $32 \times 10 \times 400 = 128.000$ segundos ($\approx 35{,}6$ horas) de execução, contra $10 \times 10 \times 400 = 40.000$ segundos ($\approx 11{,}1$ horas) da abordagem OFAT --- uma redução de aproximadamente 69\% no esforço experimental. Reconhece-se que o OFAT não captura possíveis interações entre parâmetros, limitação discutida na Seção de conclusões. Contudo, os três parâmetros da LNS governam aspectos distintos e relativamente independentes do algoritmo --- o tamanho da vizinhança (\texttt{lns\_perc}), a estratégia de seleção de candidatos (\texttt{lns\_order}) e o orçamento computacional por chamada (\texttt{lns\_time}) ---, o que reduz a probabilidade de interações fortes e torna a abordagem adequada para este contexto.

Para a seleção do melhor valor de cada parâmetro, adotou-se a seguinte hierarquia de critérios de desempate: (i)~tempo de execução, (ii)~gap de otimalidade e (iii)~número de nós restantes na árvore de Branch-and-Bound (\emph{nodes left}). Quando os critérios de maior prioridade apresentaram empate entre os cenários avaliados, o critério subsequente foi utilizado como desempate.

\section{Análise Interna dos Parâmetros da LNS} \label{sec:analise_interna}

\subsection{Impacto da Taxa de Destruição (10\%, 25\%, 50\% e 75\%)}

A primeira análise experimental investigou o impacto da taxa de destruição (\texttt{lns\_perc}) na qualidade das soluções e no desempenho computacional. Foram avaliados quatro cenários: 10\%, 25\%, 50\% e 75\% de destruição, mantendo fixos o tempo máximo do LNS (30s por chamada), a ordenação (crescente) e o tempo total de execução (400s).

\textbf{Tempo de Execução.} A Tabela~\ref{tab:lns_cand_tempo} apresenta os tempos de execução para as 10 instâncias selecionadas. A variante B\&B puro (sem LNS) obteve o menor tempo total, com \textbf{623 segundos}, contra 753 segundos do melhor cenário LNS (75\%). Esse resultado é esperado, pois o \emph{overhead} introduzido pelas chamadas ao sub-MIP do LNS aumenta o tempo total, sobretudo em instâncias que já são resolvidas de forma ótima pelo Branch-and-Bound puro. Nove das dez instâncias foram resolvidas até a otimalidade por todos os cenários, e apenas a instância input7 atingiu o limite de tempo. Embora o B\&B puro tenha sido claramente superior, as diferenças entre os quatro cenários LNS foram pequenas (753s a 756s de tempo total), não sendo suficientes para justificar a escolha de um cenário em detrimento dos demais.

% ============================================
% TABELA 1: Impacto da Taxa de Destruição - Tempo
% ============================================
\begin{table}[htbp]
\centering
\caption{Tempo de execução (segundos) para diferentes taxas de destruição do LNS. Cenário avaliado: ordenação crescente, tempo LNS = 30s, limite global = 400s.}
\label{tab:lns_cand_tempo}
\small
\begin{tabular}{@{}lrrrrrrrr@{}}
\toprule
\textbf{Instância} & \textbf{Disciplinas} & \textbf{Profs} & \textbf{B\&B} & \textbf{10\%} & \textbf{25\%} & \textbf{50\%} & \textbf{75\%} & \textbf{Ajuste} \\
\midrule
input2  & 167 & 52 & 4,04    & 17,07  & 17,66  & 16,98  & 17,04  & 50\% \\
input5  & 214 & 69 & 113,03  & 145,27 & 145,04 & 145,71 & 143,63 & 75\% \\
input6  & 79  & 24 & 1,07    & 2,47   & 2,40   & 2,33   & 2,22   & 75\% \\
input7  & 84  & 29 & 400,02  & 400,02 & 400,02 & 400,02 & 400,02 & Empate \\
input9  & 104 & 35 & 17,06   & 72,52  & 72,33  & 72,57  & 73,48  & 25\% \\
input13 & 100 & 34 & 29,13   & 60,21  & 59,56  & 59,04  & 59,43  & 50\% \\
input19 & 95  & 29 & 2,11    & 1,39   & 1,44   & 1,42   & 1,47   & 10\% \\
input20 & 104 & 37 & 0,95    & 4,69   & 4,56   & 4,59   & 5,02   & 25\% \\
input26 & 114 & 38 & 33,73   & 25,50  & 24,67  & 24,81  & 24,38  & 75\% \\
input34 & 161 & 54 & 21,92   & 26,68  & 26,08  & 26,38  & 26,41  & 25\% \\
\midrule
% \textbf{Média}  & 114 & 39 & \textbf{19,49} & 26,09 & 25,38 & 25,59 & 25,39 & \textbf{B\&B} \\
\textbf{Total}  & 1.222 & 401 & 623 & 756 & 754 & 754 & \textbf{753} & \textbf{75\%} \\
\bottomrule
\end{tabular}
\end{table}


\textbf{Gap de Otimalidade.} Em termos de gap, todos os cenários --- incluindo o B\&B puro --- apresentaram desempenho equivalente (0,60\% de gap total), com apenas a instância input7 permanecendo sem solução ótima comprovada (ver Tabela~\ref{tab:lns_cand_gap}). Esse empate impediu o uso do gap como critério de seleção.

% ============================================
% TABELA 3: Impacto da Taxa de Destruição - GAP
% ============================================
\begin{table}[htbp]
\centering
\caption{Gap de otimalidade (\%) para diferentes taxas de destruição. Combinação de parâmetros mantida fixa: ordenação crescente, tempo LNS = 30s, limite global = 400s.}
\label{tab:lns_cand_gap}
\small
\begin{tabular}{@{}lrrrrrrrr@{}}
\toprule
\textbf{Instância} & \textbf{Disciplinas} & \textbf{Profs} & \textbf{B\&B} & \textbf{10\%} & \textbf{25\%} & \textbf{50\%} & \textbf{75\%} & \textbf{Ajuste} \\
\midrule
input2  & 167 & 52 & 0,00 & 0,00 & 0,00 & 0,00 & 0,00 & Empate \\
input5  & 214 & 69 & 0,00 & 0,00 & 0,00 & 0,00 & 0,00 & Empate \\
input6  & 79  & 24 & 0,00 & 0,00 & 0,00 & 0,00 & 0,00 & Empate \\
input7  & 84  & 29 & 0,60 & 0,60 & 0,60 & 0,60 & 0,60 & Empate \\
input9  & 104 & 35 & 0,00 & 0,00 & 0,00 & 0,00 & 0,00 & Empate \\
input13 & 100 & 34 & 0,00 & 0,00 & 0,00 & 0,00 & 0,00 & Empate \\
input19 & 95  & 29 & 0,00 & 0,00 & 0,00 & 0,00 & 0,00 & Empate \\
input20 & 104 & 37 & 0,00 & 0,00 & 0,00 & 0,00 & 0,00 & Empate \\
input26 & 114 & 38 & 0,00 & 0,00 & 0,00 & 0,00 & 0,00 & Empate \\
input34 & 161 & 54 & 0,00 & 0,00 & 0,00 & 0,00 & 0,00 & Empate \\
\midrule
% \textbf{Média}  & 114 & 39 & 0,06 & 0,06 & 0,06 & 0,06 & 0,06 & Empate \\
\textbf{Total}  & 1.222 & 401 & 0,60 & 0,60 & 0,60 & 0,60 & 0,60 & \textbf{Empate} \\
\bottomrule
\end{tabular}
\end{table}


\textbf{Nós Restantes (\emph{Nodes Left}).} Uma vez que o tempo de execução não apresentou diferenças significativas entre os cenários LNS e o gap resultou em empate absoluto, o critério de desempate utilizado foi a quantidade de nós restantes na árvore de Branch-and-Bound ao final da execução, métrica que reflete o progresso em direção à prova de otimalidade. A Tabela~\ref{tab:lns_cand_nodes} revela que o cenário com 75\% de destruição alcançou o menor total de nós restantes: \textbf{58.921 nós}, comparado a 66.230 (B\&B puro), 62.628 (10\%), 60.540 (25\%) e 61.056 (50\%). Todos os cenários obtiveram resultado idêntico em 9 das 10 instâncias (resolvidas até a otimalidade), mas na instância crítica input7, a taxa de 75\% demonstrou maior capacidade de exploração e poda da árvore. A Figura~\ref{fig:nodes_cand} apresenta visualmente a comparação dos totais, evidenciando a superioridade do cenário de 75\%.

% ============================================
% TABELA 2: Impacto da Taxa de Destruição - Nodes Left
% ============================================
\begin{table}[htbp]
\centering
\caption{Número de nós restantes na árvore de Branch-and-Bound para diferentes taxas de destruição. Ajuste dos demais parâmetros: ordenação crescente, tempo LNS = 30s, limite global = 400s.}
\label{tab:lns_cand_nodes}
\small
\begin{tabular}{@{}lrrrrrrrr@{}}
\toprule
\textbf{Instância} & \textbf{Disciplinas} & \textbf{Profs} & \textbf{B\&B} & \textbf{10\%} & \textbf{25\%} & \textbf{50\%} & \textbf{75\%} & \textbf{Ajuste} \\
\midrule
input2  & 167 & 52 & 0      & 0      & 0      & 0      & 0      & -- \\
input5  & 214 & 69 & 0      & 0      & 0      & 0      & 0      & -- \\
input6  & 79  & 24 & 0      & 0      & 0      & 0      & 0      & -- \\
input7  & 84  & 29 & 66.230 & 62.628 & 60.540 & 61.056 & \textbf{58.921} & \textbf{75\%} \\
input9  & 104 & 35 & 0      & 0      & 0      & 0      & 0      & -- \\
input13 & 100 & 34 & 0      & 0      & 0      & 0      & 0      & -- \\
input19 & 95  & 29 & 0      & 0      & 0      & 0      & 0      & -- \\
input20 & 104 & 37 & 0      & 0      & 0      & 0      & 0      & -- \\
input26 & 114 & 38 & 0      & 0      & 0      & 0      & 0      & -- \\
input34 & 161 & 54 & 0      & 0      & 0      & 0      & 0      & -- \\
\midrule
% \textbf{Média}  & 114 & 39 & 6.623 & 6.263 & 6.054 & 6.106 & \textbf{5.892} & \textbf{LNS 75\%} \\
\textbf{Total}  & 1.222 & 401 & 66.230 & 62.628 & 60.540 & 61.056 & \textbf{58.921} & \textbf{75\%} \\
\bottomrule
\end{tabular}
\end{table}


\input{Auxiliares/fig_nodes_cand.tex}

\textbf{Conclusão Parcial.} A taxa de destruição de \textbf{75\%} foi selecionada para as análises subsequentes, pois apresentou a maior redução da árvore de busca na instância difícil avaliada. Para as instâncias testadas, esse resultado sugere que vizinhanças maiores podem permitir explorações mais amplas do espaço de soluções, reforçando a estratégia de intensificação adotada pela matheurística.

\subsection{Impacto do Sentido de Ordenação (Crescente vs. Decrescente)}

Com a taxa de destruição fixada em 75\%, investigou-se o impacto do sentido de ordenação dos candidatos pelo valor de \texttt{avgPreferenceWeight} (~\eqref{eq:avg_pref_weight}). Conforme descrito no Capítulo~\ref{ch:pdd}, a ordenação crescente prioriza a liberação de professores mais versáteis (maior número de preferências declaradas), enquanto a decrescente prioriza professores mais especializados (menor número de preferências).

\textbf{Tempo de Execução.} A Tabela~\ref{tab:lns_order_tempo} revela desempenho temporal praticamente idêntico entre as duas estratégias: 1.803,48 segundos (crescente) versus 1.804,48 segundos (decrescente) no tempo total, com diferença de apenas 1 segundo. O B\&B puro obteve o menor tempo total com \textbf{1.792,24 segundos}. Das 10 instâncias testadas, quatro atingiram o limite de tempo em ambas as variantes (input4, input15, input17 e input33), e as demais apresentaram variações inferiores a 1 segundo. O tempo, portanto, não permitiu diferenciar as estratégias de ordenação.

% ============================================
% TABELA 4: Impacto da Ordenação - Tempo
% ============================================
\begin{table}[htbp]
\centering
\caption{Tempo de execução (segundos) comparando ordenação crescente e decrescente. Configuração: 75\% de destruição, tempo LNS = 30s, limite global = 400s.}
\label{tab:lns_order_tempo}
\small
\begin{tabular}{@{}lrrrrrr@{}}
\toprule
\textbf{Instância} & \textbf{Disciplinas} & \textbf{Profs} & \textbf{Relax.} & \textbf{Crescente} & \textbf{Decrescente} & \textbf{Melhor} \\
\midrule
input4  & 123 & 46 & 400,05 & 400,04 & 400,05 & Crescente \\
input5  & 214 & 69 & 114,18 & 140,64 & 141,47 & Crescente \\
input12 & 115 & 37 & 5,79   & 25,70  & 25,84  & Crescente \\
input15 & 94  & 30 & 400,03 & 400,02 & 400,02 & Decrescente \\
input16 & 89  & 28 & 3,07   & 0,72   & 0,72   & Decrescente \\
input17 & 108 & 38 & 400,04 & 400,04 & 400,03 & Decrescente \\
input20 & 104 & 37 & 0,90   & 4,40   & 4,32   & Decrescente \\
input33 & 104 & 36 & 400,03 & 400,03 & 400,03 & Decrescente \\
input37 & 152 & 53 & 68,00  & 31,61  & 31,62  & Crescente \\
input38 & 62  & 22 & 0,16   & 0,28   & 0,28   & Decrescente \\
\midrule
% \textbf{Média}  & 106 & 37 & 91,09 & 86,12 & 86,55 & Crescente \\
\textbf{Total}  & 1.165 & 396 & \textbf{1.792,24} & 1.803,48 & 1.804,48 & \textbf{Relax.} \\
\bottomrule
\end{tabular}
\end{table}


\textbf{Gap de Otimalidade.} Ambas as variantes mantiveram gap total idêntico de 1,62\%, distribuído entre quatro instâncias não resolvidas de forma ótima (input4, input15, input17 e input33), com gaps individuais variando de 0,20\% a 0,73\% (ver Tabela~\ref{tab:lns_order_gap}). O empate no gap impediu, novamente, seu uso como critério de seleção.


% ============================================
% TABELA 6: Impacto da Ordenação - GAP
% ============================================
\begin{table}[htbp]
\centering
\caption{Gap de otimalidade (\%) comparando ordenação crescente e decrescente. Configuração: 75\% de destruição, tempo LNS = 30s, limite global = 400s.}
\label{tab:lns_order_gap}
\small
\begin{tabular}{@{}lrrrrrr@{}}
\toprule
\textbf{Instância} & \textbf{Disciplinas} & \textbf{Profs} & \textbf{Relax.} & \textbf{Crescente} & \textbf{Decrescente} & \textbf{Melhor} \\
\midrule
input4  & 123 & 46 & 0,20 & 0,20 & 0,20 & Empate \\
input5  & 214 & 69 & 0,00 & 0,00 & 0,00 & Empate \\
input12 & 115 & 37 & 0,00 & 0,00 & 0,00 & Empate \\
input15 & 94  & 30 & 0,73 & 0,73 & 0,73 & Empate \\
input16 & 89  & 28 & 0,00 & 0,00 & 0,00 & Empate \\
input17 & 108 & 38 & 0,20 & 0,20 & 0,20 & Empate \\
input20 & 104 & 37 & 0,00 & 0,00 & 0,00 & Empate \\
input33 & 104 & 36 & 0,49 & 0,49 & 0,49 & Empate \\
input37 & 152 & 53 & 0,00 & 0,00 & 0,00 & Empate \\
input38 & 62  & 22 & 0,00 & 0,00 & 0,00 & Empate \\
\midrule
% \textbf{Média}  & 106 & 37 & 0,16 & 0,16 & 0,16 & \textbf{Empate} \\
\textbf{Total}  & 1.165 & 396 & 1,62 & 1,62 & 1,62 & \textbf{Empate} \\
\bottomrule
\end{tabular}
\end{table}

\textbf{Nós Restantes (\emph{Nodes Left}).} A métrica de nós restantes revelou vantagem para a ordenação crescente. A Tabela~\ref{tab:lns_order_nodes} mostra que a variante crescente resultou em \textbf{131.997 nós restantes} no total, enquanto a decrescente deixou 133.674 nós --- uma redução de 1.677 nós (1,25\%). A diferença se concentra em duas instâncias críticas: input15 (61.546 vs.\ 62.452 nós) e input33 (46.243 vs.\ 47.063 nós). A Figura~\ref{fig:nodes_order} apresenta a comparação visual dos totais.

% ============================================
% TABELA 5: Impacto da Ordenação - Nodes Left
% ============================================
\begin{table}[htbp]
\centering
\caption{Número de nós restantes comparando ordenação crescente e decrescente. Configuração: 75\% de destruição, tempo LNS = 30s, limite global = 400s.}
\label{tab:lns_order_nodes}
\small
\begin{tabular}{@{}lrrrrrr@{}}
\toprule
\textbf{Instância} & \textbf{Disciplinas} & \textbf{Profs} & \textbf{Relax.} & \textbf{Crescente} & \textbf{Decrescente} & \textbf{Melhor} \\
\midrule
input4  & 123 & 46 & 19.797 & 19.541 & 19.507 & Decrescente \\
input5  & 214 & 69 & 0      & 0      & 0      & -- \\
input12 & 115 & 37 & 0      & 0      & 0      & -- \\
input15 & 94  & 30 & 64.711 & 61.546 & 62.452 & Crescente \\
input16 & 89  & 28 & 0      & 0      & 0      & -- \\
input17 & 108 & 38 & 4.673  & 4.667  & 4.652  & Decrescente \\
input20 & 104 & 37 & 0      & 0      & 0      & -- \\
input33 & 104 & 36 & 50.279 & 46.243 & 47.063 & Crescente \\
input37 & 152 & 53 & 0      & 0      & 0      & -- \\
input38 & 62  & 22 & 0      & 0      & 0      & -- \\
\midrule
% \textbf{Média}  & 106 & 37 & 13.946 & \textbf{13.200} & 13.367 & \textbf{Crescente} \\
\textbf{Total}  & 1.165 & 396 & 139.460 & \textbf{131.997} & 133.674 & \textbf{Crescente} \\
\bottomrule
\end{tabular}
\end{table}

\input{Auxiliares/fig_nodes_order.tex}

\textbf{Interpretação.} A superioridade da ordenação \textbf{crescente} sugere que, para o PDD, priorizar a realocação de professores mais versáteis (generalistas) oferece maior flexibilidade para o sub-MIP encontrar rearranjos que melhorem a solução global. Professores generalistas possuem preferências por um maior número de disciplinas, o que amplia o espaço de busca do sub-MIP e facilita a satisfação simultânea de múltiplas restrições de carga horária. A ordenação crescente foi, portanto, adotada para as análises subsequentes.

\subsection{Impacto do Tempo Máximo do LNS (10s, 50s, 100s e 200s)}

Com os parâmetros de taxa de destruição (75\%) e ordenação (crescente) definidos, investigou-se o impacto do limite de tempo do sub-MIP (\texttt{lns\_time}). Foram testados quatro valores: 10s, 50s, 100s e 200s, mantendo o limite global de 400 segundos por instância.

\textbf{Tempo de Execução.} A Tabela~\ref{tab:lns_time_tempo} apresenta resultados contraintuitivos: embora se esperasse que limites maiores de LNS aumentassem o tempo total, observou-se comportamento não monotônico. O tempo total foi de 1.629,17s (B\&B puro), 1.647,01s (10s), 1.313,60s (50s), \textbf{1.039,68s (100s)} e 1.090,00s (200s). Destaca-se a instância input15, em que os cenários de 100s e 200s resolveram o problema em $\sim$77s, contra 400s nos demais. O ajuste de 100s obteve o menor tempo total, porém a diferença para 200s (50,32s) foi pequena, não permitindo uma conclusão unívoca por este critério.

% ============================================
% TABELA 7: Impacto do Tempo Máximo LNS - Tempo
% ============================================
\begin{table}[htbp]
\centering
\caption{Tempo de execução (segundos) para diferentes limites de tempo do sub-MIP. Configuração: 75\% de destruição, ordenação crescente, limite global = 400s.}
\label{tab:lns_time_tempo}
\small
\begin{tabular}{@{}lrrrrrrrrr@{}}
\toprule
\textbf{Inst.} & \textbf{Disc.} & \textbf{Prof.} & \textbf{Relax.} & \textbf{10s} & \textbf{50s} & \textbf{100s} & \textbf{200s} & \textbf{Melhor} \\
\midrule
input6  & 79  & 24 & 1,19   & 2,40   & 2,14   & 2,15   & 2,14   & 200s \\
input14 & 99  & 28 & 14,21  & 22,31  & 68,47  & 118,44 & 168,12 & Relax. \\
input15 & 94  & 30 & 400,03 & 400,02 & 400,02 & 76,31 & 77,09 & 100s \\
input21 & 104 & 31 & 6,98   & 3,02   & 3,04   & 3,05   & 3,10   & 10s \\
input22 & 133 & 41 & 400,05 & 400,04 & 400,04 & 400,04 & 400,04 & 10s \\
input23 & 118 & 35 & 1,55   & 11,41  & 12,80  & 12,66  & 12,50  & 10s \\
input27 & 125 & 39 & 4,93   & 7,45   & 7,34   & 7,33   & 7,31   & 200s \\
input29 & 103 & 39 & 400,03 & 400,03 & 19,42  & 19,40  & 19,39  & 200s \\
input33 & 104 & 36 & 400,03 & 400,03 & 400,03 & 400,03 & 400,03 & 200s \\
input38 & 62  & 22 & 0,16   & 0,28   & 0,28   & 0,28   & 0,28   & 100s \\
\midrule
% \textbf{Média} & 104 & 33 & 162,92 & 164,77 & 162,11 & 160,27 & \textbf{159,46} & \textbf{200s} \\
\textbf{Total} & 1,021 & 325 & 1.629,17 & 1.647,01 & 1.313,60 & \textbf{1.039,68} & 1.090,00 & \textbf{100s} \\
\bottomrule
\end{tabular}
\end{table}


Essa redução de tempo com limites maiores de LNS pode ser explicada pela capacidade de encontrar soluções de melhor qualidade mais cedo: uma incumbente mais forte permite podas mais agressivas na árvore de Branch-and-Bound, acelerando a convergência. Cenários com tempo insuficiente (10s) produzem melhorias marginais que não compensam o \emph{overhead} da heurística.

\textbf{Gap de Otimalidade.} A Tabela~\ref{tab:lns_time_gap} mostra que os cenários de 100s e 200s empataram com gap total de \textbf{0,72\%}, substancialmente inferior ao obtido pelo B\&B puro e pelo LNS com 10s (ambos com 1,73\%). O destaque vai para a instância input15, resolvida até a otimalidade (gap 0,00\%) apenas pelos cenários de 100s e 200s, e para a instância input33, cujo gap caiu de 0,49\% (B\&B puro) para 0,24\% com 100s e 200s. Entretanto, o empate entre 100s e 200s impediu o uso do gap como critério final.

% ============================================
% TABELA 9: Impacto do Tempo Máximo LNS - GAP
% ============================================
\begin{table}[htbp]
\centering
\caption{Gap de otimalidade (\%) para diferentes limites de tempo do sub-MIP. Configuração: 75\% de destruição, ordenação crescente, limite global = 400s.}
\label{tab:lns_time_gap}
\small
\begin{tabular}{@{}lrrrrrrrrr@{}}
\toprule
\textbf{Inst.} & \textbf{Disc.} & \textbf{Prof.} & \textbf{Relax.} & \textbf{10s} & \textbf{50s} & \textbf{100s} & \textbf{200s} & \textbf{Melhor} \\
\midrule
input6  & 79  & 24 & 0,00 & 0,00 & 0,00 & 0,00 & 0,00 & Empate \\
input14 & 99  & 28 & 0,00 & 0,00 & 0,00 & 0,00 & 0,00 & Empate \\
input15 & 94  & 30 & 0,73 & 0,73 & 0,73 & 0,00 & 0,00 & 100s \\
input21 & 104 & 31 & 0,00 & 0,00 & 0,00 & 0,00 & 0,00 & Empate \\
input22 & 133 & 41 & 0,32 & 0,32 & 0,32 & 0,48 & 0,48 & Relax. \\
input23 & 118 & 35 & 0,00 & 0,00 & 0,00 & 0,00 & 0,00 & Empate \\
input27 & 125 & 39 & 0,00 & 0,00 & 0,00 & 0,00 & 0,00 & Empate \\
input29 & 103 & 39 & 0,19 & 0,19 & 0,00 & 0,00 & 0,00 & 50s \\
input33 & 104 & 36 & 0,49 & 0,49 & 0,49 & 0,24 & 0,24 & 100s \\
input38 & 62  & 22 & 0,00 & 0,00 & 0,00 & 0,00 & 0,00 & Empate \\
\midrule
% \textbf{Média} & 104 & 33 & 0,17 & 0,17 & 0,15 & \textbf{0,07} & \textbf{0,07} & \textbf{Empate} \\
\textbf{Total} & 1.021 & 325 & 1,73 & 1,73 & 1,54 & \textbf{0,72} & \textbf{0,72} & \textbf{Empate} \\
\bottomrule
\end{tabular}
\end{table}


\textbf{Nós Restantes (\emph{Nodes Left}).} A Tabela~\ref{tab:lns_time_nodes} foi o critério de desempate definitivo. O total de nós restantes foi: 135.986 (B\&B puro), 144.674 (10s), 130.424 (50s), 51.935 (100s) e \textbf{33.539 (200s)}. O cenário de 200s reduziu os nós em 75,3\% em relação ao B\&B puro e em 76,8\% em relação ao LNS com 10s. As instâncias input15 e input29 foram resolvidas completamente (0 nós restantes) pelos cenários de 100s e 200s, enquanto permaneceram com dezenas de milhares de nós nos cenários inferiores. Na instância input33, o cenário de 200s deixou apenas 539 nós, contra 1.524 com 100s. A Figura~\ref{fig:nodes_time} apresenta a comparação visual dos totais, evidenciando a clara superioridade do limite de 200s.


% ============================================
% TABELA 8: Impacto do Tempo Máximo LNS - Nodes Left
% ============================================
\begin{table}[htbp]
\centering
\caption{Número de nós restantes para diferentes limites de tempo do sub-MIP. Configuração: 75\% de destruição, ordenação crescente, limite global = 400s.}
\label{tab:lns_time_nodes}
\small
\begin{tabular}{@{}lrrrrrrrrr@{}}
\toprule
\textbf{Inst.} & \textbf{Disc.} & \textbf{Prof.} & \textbf{Relax.} & \textbf{10s} & \textbf{50s} & \textbf{100s} & \textbf{200s} & \textbf{Melhor} \\
\midrule
input6  & 79  & 24 & 0      & 0      & 0      & 0      & 0      & -- \\
input14 & 99  & 28 & 0      & 0      & 0      & 0      & 0      & -- \\
input15 & 94  & 30 & 54.355 & 63.056 & 59.819 & 0      & 0      & 100s \\
input21 & 104 & 31 & 0      & 0      & 0      & 0      & 0      & -- \\
input22 & 133 & 41 & 21.528 & 29.161 & 25.100 & 50.411 & 33.000 & 50s \\
input23 & 118 & 35 & 0      & 0      & 0      & 0      & 0      & -- \\
input27 & 125 & 39 & 0      & 0      & 0      & 0      & 0      & -- \\
input29 & 103 & 39 & 12.643 & 4.347  & 0      & 0      & 0      & 50s \\
input33 & 104 & 36 & 47.460 & 48.110 & 45.505 & 1.524  & 539 & 200s \\
input38 & 62  & 22 & 0      & 0      & 0      & 0      & 0      & -- \\
\midrule
% \textbf{Média} & 104 & 33 & 13.599 & 14.467 & 13.042 & 5.194 & \textbf{3.354} & \textbf{200s} \\
\textbf{Total} & 1.021 & 325 & 135.986 & 144.674 & 130.424 & 51.935 & \textbf{33.539} & \textbf{200s} \\
\bottomrule
\end{tabular}
\end{table}


\input{Auxiliares/fig_nodes_time.tex}

\textbf{Conclusão Parcial.} Entre os cenários avaliados, o limite de \textbf{200 segundos} apresentou o melhor desempenho global, reduzindo a árvore de busca e resolvendo instâncias que os demais cenários não resolveram. Para as instâncias testadas, os resultados sugerem que investir mais tempo em cada chamada do LNS pode produzir soluções de maior qualidade e acelerar a convergência do Branch-and-Bound. Nesse contexto experimental, o \emph{trade-off} entre intensidade (tempo por chamada) e frequência (número de chamadas) favoreceu a intensidade.

\subsection{Resumo da Combinação de Parâmetros Selecionada}

A Tabela~\ref{tab:config_final} resume os parâmetros selecionados ao longo da análise incremental.

\begin{table}[htbp]
\centering
\caption{Combinação final de parâmetros da heurística LNS selecionada pela análise incremental.}
\label{tab:config_final}
\small
\begin{tabular}{@{}llll@{}}
\toprule
\textbf{Parâmetro} & \textbf{Valores Testados} & \textbf{Selecionado} & \textbf{Critério Decisivo} \\
\midrule
Taxa de destruição (\texttt{lns\_perc}) & 10\%, 25\%, 50\%, 75\% & 75\% & Nós restantes \\
Ordenação (\texttt{lns\_order})         & Crescente, Decrescente  & Crescente & Nós restantes \\
Tempo do sub-MIP (\texttt{lns\_time})   & 10s, 50s, 100s, 200s   & 200s & Nós restantes \\
\bottomrule
\end{tabular}
\end{table}


Nota-se que, em todas as etapas, o tempo de execução não apresentou diferenças significativas entre os cenários LNS e o gap de otimalidade resultou em empate, sendo a métrica de nós restantes o critério de desempate definitivo. Esse padrão sugere que, para as instâncias avaliadas, o LNS atua predominantemente como acelerador da prova de otimalidade, reduzindo o esforço do Branch-and-Bound por meio de incumbentes de melhor qualidade.

\section{Comparativo de Desempenho: LNS, GRASP e Modelo Exato} \label{sec:comparativo}

Com os parâmetros da LNS calibrados (75\% de destruição, ordenação crescente e tempo de sub-MIP de 200s), realizou-se um comparativo final entre quatro variantes: (i)~\textbf{B\&B puro} (sem heurísticas primais customizadas), (ii)~\textbf{LNS} (apenas a heurística LNS proposta), (iii)~\textbf{GRASP} (apenas a heurística GRASP, com \texttt{max\_iter=1}, $\alpha=0{,}8$ e sem busca local \cite{jhonatan_grasp_dpd}) e (iv)~\textbf{GRASP|LNS} (ambas as heurísticas ativadas simultaneamente). As quatro variantes foram executadas sobre as mesmas 43 instâncias: 42 sintéticas e uma real. Nas tabelas detalhadas, a instância real é identificada explicitamente e destacada em cinza; seus resultados também integram os totais e as médias gerais.

\subsection{Análise do Tempo de Execução}

A Tabela~\ref{tab:geral_tempo} apresenta os tempos de execução por instância. Considerando o conjunto completo, o \textbf{B\&B puro} obteve o menor tempo total, com \textbf{5.232,81 segundos}, seguido por LNS (6.102,82s), GRASP|LNS (6.176,44s) e GRASP (6.989,55s). Em relação ao B\&B puro, os tempos totais foram 16,6\% maiores para o LNS, 18,0\% maiores para o GRASP|LNS e 33,6\% maiores para o GRASP. O resultado evidencia que o custo das heurísticas não foi compensado por reduções sistemáticas no tempo de resolução ao longo das 43 instâncias.

\begingroup
\setlength{\tabcolsep}{3pt}
\scriptsize
\begin{longtable}{@{}lrrrrrrl@{}}
\caption{Tempo de execução (segundos) nas 42 instâncias sintéticas e na instância real, destacada em cinza.}\label{tab:geral_tempo}\\
\toprule
\textbf{Instância} & \textbf{Disc.} & \textbf{Profs} & \textbf{B\&B} & \textbf{LNS} & \textbf{GRASP} & \textbf{GRASP|LNS} & \textbf{Melhor} \\
\midrule
\endfirsthead
\multicolumn{8}{c}{\tablename\ \thetable{} -- continuação}\\
\toprule
\textbf{Instância} & \textbf{Disc.} & \textbf{Profs} & \textbf{B\&B} & \textbf{LNS} & \textbf{GRASP} & \textbf{GRASP|LNS} & \textbf{Melhor} \\
\midrule
\endhead
\midrule
\multicolumn{8}{r}{Continua na próxima página.}\\
\endfoot
\bottomrule
\endlastfoot
input1 & 194 & 51 & 13,87 & 14,04 & 55,33 & 32,23 & B\&B \\
input2 & 167 & 52 & 5,07 & 5,97 & 19,07 & 9,25 & B\&B \\
input3 & 180 & 50 & 3,15 & 3,12 & 11,20 & 11,26 & LNS \\
input4 & 123 & 46 & 11,45 & 208,94 & 38,47 & 228,73 & B\&B \\
input5 & 214 & 69 & 33,99 & 212,68 & 122,70 & 253,37 & B\&B \\
input6 & 79 & 24 & 0,73 & 2,00 & 1,79 & 0,49 & GRASP{|}LNS \\
input7 & 84 & 29 & 400,02 & 400,02 & 400,02 & 400,02 & Empate \\
input8 & 166 & 46 & 400,06 & 209,41 & 400,07 & 231,67 & LNS \\
input9 & 104 & 35 & 15,30 & 55,09 & 50,07 & 54,58 & B\&B \\
input10 & 113 & 31 & 1,11 & 1,22 & 3,39 & 0,73 & GRASP{|}LNS \\
input11 & 100 & 31 & 3,41 & 1,14 & 0,95 & 0,57 & GRASP{|}LNS \\
input12 & 115 & 37 & 4,76 & 6,17 & 10,93 & 1,76 & GRASP{|}LNS \\
input13 & 100 & 34 & 98,87 & 242,55 & 252,49 & 45,39 & GRASP{|}LNS \\
input14 & 99 & 28 & 54,35 & 66,69 & 131,38 & 139,80 & B\&B \\
input15 & 94 & 30 & 45,62 & 113,87 & 400,03 & 58,63 & B\&B \\
input16 & 89 & 28 & 0,24 & 0,34 & 0,46 & 0,37 & B\&B \\
input17 & 108 & 38 & 400,04 & 400,04 & 400,04 & 400,04 & Empate \\
input18 & 121 & 38 & 20,21 & 3,00 & 7,96 & 4,32 & LNS \\
input19 & 95 & 29 & 1,52 & 1,25 & 2,40 & 0,72 & GRASP{|}LNS \\
input20 & 104 & 37 & 9,56 & 19,69 & 20,31 & 19,71 & B\&B \\
input21 & 104 & 31 & 8,24 & 4,84 & 18,14 & 7,58 & LNS \\
input22 & 133 & 41 & 400,05 & 400,05 & 400,05 & 400,05 & Empate \\
input23 & 118 & 35 & 1,32 & 1,78 & 2,87 & 1,07 & GRASP{|}LNS \\
input24 & 117 & 33 & 40,76 & 400,03 & 400,04 & 400,03 & B\&B \\
input25 & 107 & 34 & 5,43 & 39,20 & 9,78 & 38,95 & B\&B \\
input26 & 114 & 38 & 400,04 & 400,04 & 400,04 & 400,04 & Empate \\
input27 & 125 & 39 & 6,44 & 8,23 & 19,44 & 11,18 & B\&B \\
input28 & 114 & 39 & 78,83 & 400,04 & 400,04 & 400,04 & B\&B \\
input29 & 103 & 39 & 4,19 & 9,01 & 8,60 & 50,69 & B\&B \\
input30 & 109 & 36 & 10,68 & 2,77 & 19,29 & 3,20 & LNS \\
input31 & 113 & 40 & 1,55 & 1,48 & 2,86 & 1,87 & LNS \\
input32 & 97 & 36 & 400,03 & 400,03 & 400,03 & 400,03 & Empate \\
input33 & 104 & 36 & 400,03 & 400,03 & 400,03 & 400,04 & Empate \\
input34 & 161 & 54 & 4,33 & 9,28 & 12,56 & 16,99 & B\&B \\
input35 & 169 & 53 & 400,08 & 400,08 & 400,08 & 400,09 & Empate \\
input36 & 185 & 60 & 5,10 & 6,64 & 15,71 & 14,78 & B\&B \\
input37 & 152 & 53 & 290,53 & 40,14 & 400,07 & 113,44 & LNS \\
input38 & 62 & 22 & 0,18 & 0,45 & 0,65 & 0,41 & B\&B \\
input39 & 130 & 43 & 53,71 & 12,11 & 146,95 & 15,22 & LNS \\
input40 & 126 & 42 & 400,05 & 400,05 & 400,05 & 400,05 & Empate \\
input41 & 137 & 42 & 400,05 & 400,05 & 400,05 & 400,05 & Empate \\
input42 & 142 & 35 & 1,55 & 5,89 & 2,98 & 6,91 & B\&B \\
\rowcolor{realrow} \textbf{entrada\_real (Real)} & 224 & 61 & 396,31 & 393,38 & 400,14 & 400,12 & LNS \\
\midrule
Total sintético & -- & -- & 4.836,50 & 5.709,44 & 6.589,40 & 5.776,32 & B\&B \\
\textbf{Total geral} & -- & -- & \textbf{5.232,81} & \textbf{6.102,82} & \textbf{6.989,55} & \textbf{6.176,44} & \textbf{B\&B} \\
\end{longtable}
\endgroup


Na instância real, o LNS apresentou o menor tempo, com 393,38s, seguido pelo B\&B puro, com 396,31s, uma diferença de 2,93s (0,74\%). GRASP e GRASP|LNS atingiram aproximadamente 400s. Portanto, embora o B\&B puro tenha sido mais eficiente no agregado, o LNS obteve pequena vantagem temporal no caso real avaliado. A Figura~\ref{fig:geral_tempo} sintetiza os tempos totais.

\input{Auxiliares/fig_geral_tempo.tex}

\subsection{Análise do Gap de Otimalidade}

A Tabela~\ref{tab:geral_gap} apresenta o gap de otimalidade em cada instância. Como gaps de instâncias distintas não constituem uma única taxa acumulável, a comparação agregada utiliza o \textbf{gap médio percentual}, acompanhado da quantidade de instâncias resolvidas até a otimalidade. O B\&B puro obteve o menor gap médio, 0,3473\%, seguido por LNS (0,4682\%), GRASP|LNS (0,5420\%) e GRASP (0,7099\%). B\&B puro e LNS resolveram 32 das 43 instâncias até a otimalidade, enquanto GRASP|LNS resolveu 31 e GRASP, 28.

\begingroup
\setlength{\tabcolsep}{3pt}
\scriptsize
\begin{longtable}{@{}lrrrrrrl@{}}
\caption{Gap de otimalidade (\%) nas 42 instâncias sintéticas e na instância real, destacada em cinza.}\label{tab:geral_gap}\\
\toprule
\textbf{Instância} & \textbf{Disc.} & \textbf{Profs} & \textbf{B\&B} & \textbf{LNS} & \textbf{GRASP} & \textbf{GRASP|LNS} & \textbf{Melhor} \\
\midrule
\endfirsthead
\multicolumn{8}{c}{\tablename\ \thetable{} -- continuação}\\
\toprule
\textbf{Instância} & \textbf{Disc.} & \textbf{Profs} & \textbf{B\&B} & \textbf{LNS} & \textbf{GRASP} & \textbf{GRASP|LNS} & \textbf{Melhor} \\
\midrule
\endhead
\midrule
\multicolumn{8}{r}{Continua na próxima página.}\\
\endfoot
\bottomrule
\endlastfoot
input1 & 194 & 51 & 0,0000 & 0,0000 & 0,0000 & 0,0000 & Empate \\
input2 & 167 & 52 & 0,0000 & 0,0000 & 0,0000 & 0,0000 & Empate \\
input3 & 180 & 50 & 0,0000 & 0,0000 & 0,0000 & 0,0000 & Empate \\
input4 & 123 & 46 & 0,0000 & 0,0000 & 0,0000 & 0,0000 & Empate \\
input5 & 214 & 69 & 0,0000 & 0,0000 & 0,0000 & 0,0000 & Empate \\
input6 & 79 & 24 & 0,0000 & 0,0000 & 0,0000 & 0,0000 & Empate \\
input7 & 84 & 29 & 0,2994 & 1,8237 & 0,2994 & 1,8237 & Empate \\
input8 & 166 & 46 & 0,1592 & 0,0000 & 0,1592 & 0,0000 & Empate \\
input9 & 104 & 35 & 0,0000 & 0,0000 & 0,0000 & 0,0000 & Empate \\
input10 & 113 & 31 & 0,0000 & 0,0000 & 0,0000 & 0,0000 & Empate \\
input11 & 100 & 31 & 0,0000 & 0,0000 & 0,0000 & 0,0000 & Empate \\
input12 & 115 & 37 & 0,0000 & 0,0000 & 0,0000 & 0,0000 & Empate \\
input13 & 100 & 34 & 0,0000 & 0,0000 & 0,0000 & 0,0000 & Empate \\
input14 & 99 & 28 & 0,0000 & 0,0000 & 0,0000 & 0,0000 & Empate \\
input15 & 94 & 30 & 0,0000 & 0,0000 & 0,2410 & 0,0000 & Empate \\
input16 & 89 & 28 & 0,0000 & 0,0000 & 0,0000 & 0,0000 & Empate \\
input17 & 108 & 38 & 0,3893 & 1,9923 & 2,1963 & 1,9923 & B\&B \\
input18 & 121 & 38 & 0,0000 & 0,0000 & 0,0000 & 0,0000 & Empate \\
input19 & 95 & 29 & 0,0000 & 0,0000 & 0,0000 & 0,0000 & Empate \\
input20 & 104 & 37 & 0,0000 & 0,0000 & 0,0000 & 0,0000 & Empate \\
input21 & 104 & 31 & 0,0000 & 0,0000 & 0,0000 & 0,0000 & Empate \\
input22 & 133 & 41 & 4,8414 & 4,8414 & 4,8414 & 4,8414 & Empate \\
input23 & 118 & 35 & 0,0000 & 0,0000 & 0,0000 & 0,0000 & Empate \\
input24 & 117 & 33 & 3,6163 & 0,2008 & 3,3769 & 0,2008 & Empate \\
input25 & 107 & 34 & 0,0000 & 0,0000 & 0,0000 & 0,0000 & Empate \\
input26 & 114 & 38 & 0,2079 & 5,2402 & 5,2402 & 5,2402 & B\&B \\
input27 & 125 & 39 & 0,0000 & 0,0000 & 0,0000 & 0,0000 & Empate \\
input28 & 114 & 39 & 0,0000 & 0,1805 & 0,1805 & 0,1805 & B\&B \\
input29 & 103 & 39 & 0,0000 & 0,0000 & 0,0000 & 0,0000 & Empate \\
input30 & 109 & 36 & 0,0000 & 0,0000 & 0,0000 & 0,0000 & Empate \\
input31 & 113 & 40 & 0,0000 & 0,0000 & 0,0000 & 0,0000 & Empate \\
input32 & 97 & 36 & 0,2597 & 0,2597 & 1,8470 & 0,2597 & Empate \\
input33 & 104 & 36 & 0,2427 & 0,4866 & 0,7317 & 0,2427 & Empate \\
input34 & 161 & 54 & 0,0000 & 0,0000 & 0,0000 & 0,0000 & Empate \\
input35 & 169 & 53 & 0,7849 & 0,1560 & 1,1024 & 0,1560 & Empate \\
input36 & 185 & 60 & 0,0000 & 0,0000 & 0,0000 & 0,0000 & Empate \\
input37 & 152 & 53 & 0,0000 & 0,0000 & 2,0043 & 0,0000 & Empate \\
input38 & 62 & 22 & 0,0000 & 0,0000 & 0,0000 & 0,0000 & Empate \\
input39 & 130 & 43 & 0,0000 & 0,0000 & 0,0000 & 0,0000 & Empate \\
input40 & 126 & 42 & 0,2193 & 0,2193 & 0,2193 & 0,2193 & Empate \\
input41 & 137 & 42 & 3,9139 & 4,7337 & 4,7337 & 4,7337 & B\&B \\
input42 & 142 & 35 & 0,0000 & 0,0000 & 0,0000 & 0,0000 & Empate \\
\rowcolor{realrow} \textbf{entrada\_real (Real)} & 224 & 61 & 0,0000 & 0,0000 & 3,3528 & 3,4147 & Empate \\
\midrule
Média sintética & -- & -- & 0,3556 & 0,4794 & 0,6470 & 0,4736 & B\&B \\
\textbf{Média geral} & -- & -- & \textbf{0,3473} & \textbf{0,4682} & \textbf{0,7099} & \textbf{0,5420} & \textbf{B\&B} \\
\end{longtable}
\endgroup


\input{Auxiliares/fig_geral_gap.tex}

Na instância real, B\&B puro e LNS alcançaram gap de 0,0000\%, ao passo que GRASP encerrou com 3,3528\% e GRASP|LNS com 3,4147\%. Esse resultado reforça que a presença das duas heurísticas em conjunto não assegura melhoria da qualidade da incumbente ou da prova de otimalidade. A Figura~\ref{fig:geral_gap} apresenta os gaps médios das 43 instâncias.

\subsection{Análise dos Nós Restantes}

A Tabela~\ref{tab:geral_nodes} mostra que a variante \textbf{GRASP|LNS} obteve o menor total de nós restantes, com \textbf{187.673 nós}, contra 416.125 do B\&B puro, 462.967 do LNS e 467.336 do GRASP. A redução da variante híbrida foi de 54,9\% em relação ao B\&B puro, 59,5\% em relação ao LNS e 59,8\% em relação ao GRASP. Embora não tenha apresentado o menor gap médio nem o menor tempo total, a combinação GRASP|LNS avançou mais na exploração e poda das árvores que permaneceram abertas.

\begingroup
\setlength{\tabcolsep}{3pt}
\scriptsize
\begin{longtable}{@{}lrrrrrrl@{}}
\caption{Número de nós restantes nas 42 instâncias sintéticas e na instância real, destacada em cinza.}\label{tab:geral_nodes}\\
\toprule
\textbf{Instância} & \textbf{Disc.} & \textbf{Profs} & \textbf{B\&B} & \textbf{LNS} & \textbf{GRASP} & \textbf{GRASP|LNS} & \textbf{Melhor} \\
\midrule
\endfirsthead
\multicolumn{8}{c}{\tablename\ \thetable{} -- continuação}\\
\toprule
\textbf{Instância} & \textbf{Disc.} & \textbf{Profs} & \textbf{B\&B} & \textbf{LNS} & \textbf{GRASP} & \textbf{GRASP|LNS} & \textbf{Melhor} \\
\midrule
\endhead
\midrule
\multicolumn{8}{r}{Continua na próxima página.}\\
\endfoot
\bottomrule
\endlastfoot
input1 & 194 & 51 & 0 & 0 & 0 & 0 & Empate \\
input2 & 167 & 52 & 0 & 0 & 0 & 0 & Empate \\
input3 & 180 & 50 & 0 & 0 & 0 & 0 & Empate \\
input4 & 123 & 46 & 0 & 0 & 0 & 0 & Empate \\
input5 & 214 & 69 & 0 & 0 & 0 & 0 & Empate \\
input6 & 79 & 24 & 0 & 0 & 0 & 0 & Empate \\
input7 & 84 & 29 & 1.405 & 69.086 & 2.027 & 37.759 & B\&B \\
input8 & 166 & 46 & 7.283 & 0 & 4.552 & 0 & Empate \\
input9 & 104 & 35 & 0 & 0 & 0 & 0 & Empate \\
input10 & 113 & 31 & 0 & 0 & 0 & 0 & Empate \\
input11 & 100 & 31 & 0 & 0 & 0 & 0 & Empate \\
input12 & 115 & 37 & 0 & 0 & 0 & 0 & Empate \\
input13 & 100 & 34 & 0 & 0 & 0 & 0 & Empate \\
input14 & 99 & 28 & 0 & 0 & 0 & 0 & Empate \\
input15 & 94 & 30 & 0 & 0 & 2.843 & 0 & Empate \\
input16 & 89 & 28 & 0 & 0 & 0 & 0 & Empate \\
input17 & 108 & 38 & 34.042 & 64.524 & 52.649 & 31.854 & GRASP{|}LNS \\
input18 & 121 & 38 & 0 & 0 & 0 & 0 & Empate \\
input19 & 95 & 29 & 0 & 0 & 0 & 0 & Empate \\
input20 & 104 & 37 & 0 & 0 & 0 & 0 & Empate \\
input21 & 104 & 31 & 0 & 0 & 0 & 0 & Empate \\
input22 & 133 & 41 & 152.910 & 91.398 & 56.289 & 28.601 & GRASP{|}LNS \\
input23 & 118 & 35 & 0 & 0 & 0 & 0 & Empate \\
input24 & 117 & 33 & 14.077 & 3.083 & 66.822 & 2.147 & GRASP{|}LNS \\
input25 & 107 & 34 & 0 & 0 & 0 & 0 & Empate \\
input26 & 114 & 38 & 3.853 & 79.390 & 64.343 & 33.110 & B\&B \\
input27 & 125 & 39 & 0 & 0 & 0 & 0 & Empate \\
input28 & 114 & 39 & 0 & 6.138 & 3.895 & 8.344 & B\&B \\
input29 & 103 & 39 & 0 & 0 & 0 & 0 & Empate \\
input30 & 109 & 36 & 0 & 0 & 0 & 0 & Empate \\
input31 & 113 & 40 & 0 & 0 & 0 & 0 & Empate \\
input32 & 97 & 36 & 2.327 & 1.797 & 74.630 & 3.380 & LNS \\
input33 & 104 & 36 & 2.514 & 36.423 & 25.783 & 3.287 & B\&B \\
input34 & 161 & 54 & 0 & 0 & 0 & 0 & Empate \\
input35 & 169 & 53 & 21.909 & 5.593 & 12.876 & 4.965 & GRASP{|}LNS \\
input36 & 185 & 60 & 0 & 0 & 0 & 0 & Empate \\
input37 & 152 & 53 & 0 & 0 & 29.896 & 0 & Empate \\
input38 & 62 & 22 & 0 & 0 & 0 & 0 & Empate \\
input39 & 130 & 43 & 0 & 0 & 0 & 0 & Empate \\
input40 & 126 & 42 & 1.453 & 1.267 & 1.118 & 775 & GRASP{|}LNS \\
input41 & 137 & 42 & 174.352 & 104.268 & 57.079 & 27.989 & GRASP{|}LNS \\
input42 & 142 & 35 & 0 & 0 & 0 & 0 & Empate \\
\rowcolor{realrow} \textbf{entrada\_real (Real)} & 224 & 61 & 0 & 0 & 12.534 & 5.462 & Empate \\
\midrule
Total sintético & -- & -- & 416.125 & 462.967 & 454.802 & 182.211 & GRASP{|}LNS \\
\textbf{Total geral} & -- & -- & \textbf{416.125} & \textbf{462.967} & \textbf{467.336} & \textbf{187.673} & \textbf{GRASP{|}LNS} \\
\end{longtable}
\endgroup


\input{Auxiliares/fig_geral_nodes.tex}

Na instância real, B\&B puro e LNS encerraram a árvore de busca sem nós restantes. GRASP deixou 12.534 nós e GRASP|LNS, 5.462, mostrando que a combinação com LNS reduziu a árvore remanescente do GRASP, mas não foi suficiente para provar a otimalidade dentro do limite adotado.

\subsection{Análise das Execuções da LNS}

A Tabela~\ref{tab:geral_lns_executions} apresenta o número de execuções efetivas da LNS. Essa métrica é aplicável apenas às variantes LNS e GRASP|LNS; por esse motivo, B\&B puro e GRASP são indicados por ``--''. Nas 42 instâncias sintéticas, foram registradas 88 execuções no LNS e 80 no GRASP|LNS. Na instância real, ocorreram 4 execuções no LNS e 10 no GRASP|LNS, resultando em totais gerais de 92 e 90, respectivamente.

\begingroup
\setlength{\tabcolsep}{3pt}
\scriptsize
\begin{longtable}{@{}lrrrrrr@{}}
\caption{Número de execuções efetivas da LNS nas 42 instâncias sintéticas e na instância real. O símbolo -- indica que a abordagem não utiliza LNS; a instância real aparece destacada.}\label{tab:geral_lns_executions}\\
\toprule
\textbf{Instância} & \textbf{Disc.} & \textbf{Profs} & \textbf{B\&B} & \textbf{LNS} & \textbf{GRASP} & \textbf{GRASP|LNS} \\
\midrule
\endfirsthead
\multicolumn{7}{c}{\tablename\ \thetable{} -- continuação}\\
\toprule
\textbf{Instância} & \textbf{Disc.} & \textbf{Profs} & \textbf{B\&B} & \textbf{LNS} & \textbf{GRASP} & \textbf{GRASP|LNS} \\
\midrule
\endhead
\midrule
\multicolumn{7}{r}{Continua na próxima página.}\\
\endfoot
\bottomrule
\endlastfoot
input1 & 194 & 51 & -- & 2 & -- & 2 \\
input2 & 167 & 52 & -- & 1 & -- & 1 \\
input3 & 180 & 50 & -- & 0 & -- & 0 \\
input4 & 123 & 46 & -- & 2 & -- & 2 \\
input5 & 214 & 69 & -- & 2 & -- & 2 \\
input6 & 79 & 24 & -- & 1 & -- & 1 \\
input7 & 84 & 29 & -- & 7 & -- & 7 \\
input8 & 166 & 46 & -- & 1 & -- & 1 \\
input9 & 104 & 35 & -- & 1 & -- & 2 \\
input10 & 113 & 31 & -- & 1 & -- & 1 \\
input11 & 100 & 31 & -- & 1 & -- & 1 \\
input12 & 115 & 37 & -- & 2 & -- & 1 \\
input13 & 100 & 34 & -- & 6 & -- & 2 \\
input14 & 99 & 28 & -- & 1 & -- & 1 \\
input15 & 94 & 30 & -- & 1 & -- & 2 \\
input16 & 89 & 28 & -- & 2 & -- & 1 \\
input17 & 108 & 38 & -- & 2 & -- & 2 \\
input18 & 121 & 38 & -- & 1 & -- & 1 \\
input19 & 95 & 29 & -- & 2 & -- & 1 \\
input20 & 104 & 37 & -- & 1 & -- & 1 \\
input21 & 104 & 31 & -- & 2 & -- & 2 \\
input22 & 133 & 41 & -- & 1 & -- & 1 \\
input23 & 118 & 35 & -- & 1 & -- & 1 \\
input24 & 117 & 33 & -- & 3 & -- & 3 \\
input25 & 107 & 34 & -- & 1 & -- & 1 \\
input26 & 114 & 38 & -- & 1 & -- & 1 \\
input27 & 125 & 39 & -- & 1 & -- & 1 \\
input28 & 114 & 39 & -- & 3 & -- & 2 \\
input29 & 103 & 39 & -- & 1 & -- & 1 \\
input30 & 109 & 36 & -- & 1 & -- & 1 \\
input31 & 113 & 40 & -- & 1 & -- & 1 \\
input32 & 97 & 36 & -- & 5 & -- & 3 \\
input33 & 104 & 36 & -- & 2 & -- & 2 \\
input34 & 161 & 54 & -- & 2 & -- & 2 \\
input35 & 169 & 53 & -- & 2 & -- & 2 \\
input36 & 185 & 60 & -- & 2 & -- & 2 \\
input37 & 152 & 53 & -- & 2 & -- & 2 \\
input38 & 62 & 22 & -- & 1 & -- & 1 \\
input39 & 130 & 43 & -- & 2 & -- & 2 \\
input40 & 126 & 42 & -- & 15 & -- & 15 \\
input41 & 137 & 42 & -- & 1 & -- & 1 \\
input42 & 142 & 35 & -- & 1 & -- & 1 \\
\rowcolor{realrow} \textbf{entrada\_real (Real)} & 224 & 61 & -- & 4 & -- & 10 \\
\midrule
Subtotal sintético & -- & -- & -- & 88 & -- & 80 \\
\textbf{Total geral} & -- & -- & -- & \textbf{92} & -- & \textbf{90} \\
\end{longtable}
\endgroup


O número de execuções representa a frequência com que o mecanismo de reotimização foi efetivamente acionado, e não uma medida direta de qualidade. Na instância real, por exemplo, o GRASP|LNS executou a LNS mais vezes que o LNS isolado, mas terminou com gap positivo; portanto, mais execuções não implicaram necessariamente melhor convergência.

\subsection{Discussão dos Resultados}

O comparativo completo não identifica uma abordagem dominante em todas as métricas. A Tabela~\ref{tab:resumo_comparativo} sintetiza os indicadores das 43 instâncias, incluindo a instância real.

\begin{table}[htbp]
\centering
\caption{Resumo comparativo das quatro variantes nas 42 instâncias sintéticas e na instância real.}
\label{tab:resumo_comparativo}
\small
\begin{tabular}{@{}lrrrr@{}}
\toprule
\textbf{Métrica} & \textbf{B\&B puro} & \textbf{LNS} & \textbf{GRASP} & \textbf{GRASP|LNS} \\
\midrule
Tempo total (s)        & \textbf{5.232,81} & 6.102,82 & 6.989,55 & 6.176,44 \\
Tempo médio (s)        & \textbf{121,69} & 141,93 & 162,55 & 143,64 \\
Gap médio (\%)         & \textbf{0,3473} & 0,4682 & 0,7099 & 0,5420 \\
Instâncias ótimas      & \textbf{32/43} & \textbf{32/43} & 28/43 & 31/43 \\
Nós restantes (total)  & 416.125 & 462.967 & 467.336 & \textbf{187.673} \\
Execuções LNS           & -- & 92 & -- & 90 \\
\bottomrule
\end{tabular}
\end{table}


A análise consolidada permite destacar quatro observações:

\begin{enumerate}
  \item \textbf{O B\&B puro foi mais eficiente no agregado.} A abordagem exata apresentou o menor tempo total e o menor gap médio, além de empatar com o LNS no maior número de instâncias ótimas. Para o conjunto completo, o custo adicional das heurísticas não produziu ganho temporal global.

  \item \textbf{O GRASP|LNS promoveu maior redução das árvores abertas.} A variante híbrida deixou menos da metade dos nós restantes do B\&B puro. Esse avanço, contudo, não se converteu no menor gap médio nem no menor tempo total, demonstrando que o número de nós restantes deve ser interpretado em conjunto com as demais métricas.

  \item \textbf{O LNS apresentou resultado competitivo na instância real.} LNS e B\&B puro provaram a otimalidade e encerraram a árvore, com vantagem de 2,93s para o LNS. Trata-se de evidência favorável à matheurística no caso real estudado, mas insuficiente para generalização a outros contextos institucionais.

  \item \textbf{O GRASP isolado apresentou o desempenho agregado menos favorável.} A abordagem registrou o maior tempo total, o maior gap médio, o menor número de instâncias ótimas e o maior total de nós restantes. A combinação com LNS reduziu substancialmente os nós remanescentes, mas não superou B\&B puro ou LNS em todas as métricas.
\end{enumerate}
